% MATA-LL User's Manual
% Compile with: lualatex usermanual.tex
\documentclass[11pt,a4paper]{article}

\usepackage{fontspec}
\usepackage{unicode-math}
\setmainfont{Latin Modern Roman}
\setmonofont{Latin Modern Mono}[Scale=0.88]
\setmathfont{Latin Modern Math}

\usepackage{geometry}
\geometry{margin=2.5cm}

\usepackage{hyperref}
\hypersetup{colorlinks=true, linkcolor=blue!60!black, urlcolor=blue!60!black}

\usepackage{listings}
\usepackage{xcolor}
\usepackage{booktabs}
\usepackage{enumitem}
\usepackage{titlesec}
\usepackage{framed}

\titleformat{\section}{\Large\bfseries}{\thesection}{1em}{}
\titleformat{\subsection}{\large\bfseries}{\thesubsection}{1em}{}
\titleformat{\subsubsection}{\normalsize\bfseries}{\thesubsubsection}{1em}{}

\definecolor{codebg}{HTML}{F5F5F5}
\definecolor{keyword}{HTML}{7B3294}
\definecolor{strcolor}{HTML}{2E7D32}
\definecolor{comment}{HTML}{757575}
\definecolor{notebg}{HTML}{FFF8E1}

\lstdefinelanguage{MLL}{
  keywords={data,where,let,in,case,of,if,then,else,do,import,module,
            class,instance,deriving,newtype,type,export,infixl,infixr,infix,
            forall,hiding,qualified,as},
  sensitive=true,
  morecomment=[l]{--},
  morecomment=[n]{\{-}{-\}},
  morestring=[b]",
}

\lstset{
  language=MLL,
  basicstyle=\ttfamily\small,
  backgroundcolor=\color{codebg},
  frame=single,
  framerule=0pt,
  rulecolor=\color{codebg},
  keywordstyle=\color{keyword}\bfseries,
  stringstyle=\color{strcolor},
  commentstyle=\color{comment}\itshape,
  breaklines=true,
  showstringspaces=false,
  tabsize=2,
  xleftmargin=4pt,
  xrightmargin=4pt,
  aboveskip=8pt,
  belowskip=8pt,
}

\newenvironment{note}{%
  \def\FrameCommand{\fboxsep=3mm \colorbox{notebg}}%
  \MakeFramed{\advance\hsize-\width \FrameRestore}%
  \noindent\textbf{Note:} }%
  {\endMakeFramed}

\title{\textbf{MATA-LL User's Manual}\\[6pt]
  \large A Haskell-Inspired Language Targeting Lua}
\author{MATA-LL Project}
\date{June 2026}

\begin{document}
\maketitle
\tableofcontents
\clearpage

%======================================================================
\section{Getting Started}

\subsection{What is MATA-LL?}

MATA-LL (Modest Attempt at Typesystem Augmenting the Lua Language) is a
statically-typed, non-strict functional language that compiles to standalone
Lua files.  If you know Haskell, you will feel at home immediately.  The
generated Lua requires no runtime library and runs on Lua~5.4 or LuaJIT.

\subsection{Installation}

Install from source with Cargo:

\begin{lstlisting}[language={}]
cargo install --path .
\end{lstlisting}

This builds an optimized binary and places it in \texttt{\textasciitilde/.cargo/bin/},
which is normally already on your \texttt{PATH}.

\subsection{Hello World}

Create \texttt{hello.mll}:

\begin{lstlisting}
main :: IO ()
main = putStrLn "Hello, world!"
\end{lstlisting}

Compile and run:

\begin{lstlisting}[language={}]
mll hello.mll --run
\end{lstlisting}

Or emit a Lua file:

\begin{lstlisting}[language={}]
mll hello.mll --emit-lua      # writes hello.lua
lua5.4 hello.lua               # runs standalone
\end{lstlisting}

\subsection{Command-Line Usage}

\begin{tabular}{@{}ll@{}}
\toprule
\textbf{Command} & \textbf{Effect} \\
\midrule
\texttt{mll file.mll}            & Compile, print Lua to stdout \\
\texttt{mll file.mll -{}-run}     & Compile and execute immediately \\
\texttt{mll file.mll -{}-emit-lua} & Write \texttt{file.lua} \\
\texttt{mll file.mll -L ./lib}   & Add library search path \\
\bottomrule
\end{tabular}

\subsection{The REPL}

Start with \texttt{mll-repl}:

\begin{lstlisting}[language={}]
mll> let x = 42
mll> x * 2 + 1
85
mll> :t map
map :: (a -> b) -> [a] -> [b]
\end{lstlisting}

REPL commands: \texttt{/clear}, \texttt{/decls}, \texttt{/lua EXPR},
\texttt{/source}, \texttt{/drop}, \texttt{/help}, \texttt{/quit}.
Multi-line input: end a line with \texttt{\textbackslash}.

%======================================================================
\section{Basic Syntax}

\subsection{Comments}

\begin{lstlisting}
-- Single-line comment
{- Block comment
   can span lines
   {- and nest -}
-}
\end{lstlisting}

\subsection{Layout Rule}

MATA-LL uses indentation-sensitive layout, like Haskell or Python.
Continuation lines must be indented past the start column of their
definition:

\begin{lstlisting}
longExpression :: Integer
longExpression = 1 + 2
    + 3 + 4    -- continuation: indented past 'l'
\end{lstlisting}

\subsection{Type Signatures}

All top-level definitions require explicit type signatures.  Local bindings
(\texttt{let}, \texttt{where}) are inferred automatically.

\begin{lstlisting}
double :: Integer -> Integer
double x = x * 2
\end{lstlisting}

%======================================================================
\section{Types and Data}

\subsection{Primitive Types}

\begin{tabular}{@{}lll@{}}
\toprule
\textbf{Type} & \textbf{Lua type} & \textbf{Examples} \\
\midrule
\texttt{Integer} & integer     & \texttt{42}, \texttt{-7}, \texttt{0} \\
\texttt{Number}  & float       & \texttt{3.14}, \texttt{-0.5} \\
\texttt{String}  & string      & \texttt{"hello\textbackslash n"} \\
\texttt{Bool}    & boolean     & \texttt{True}, \texttt{False} \\
\bottomrule
\end{tabular}

\begin{note}
There is no \texttt{Char} type.  Strings are native Lua strings, not
\texttt{[Char]}.  Use \texttt{<>} for string concatenation (not \texttt{++}).
\end{note}

\subsection{Lists}

Lists are cons-cell linked lists (not Lua arrays):

\begin{lstlisting}
nums :: [Integer]
nums = [1, 2, 3, 4, 5]

-- Equivalent to:
nums' :: [Integer]
nums' = 1 : 2 : 3 : 4 : 5 : []
\end{lstlisting}

\subsection{Tuples}

\begin{lstlisting}
pair :: (String, Integer)
pair = ("age", 30)

triple :: (Bool, Integer, String)
triple = (True, 42, "hello")
\end{lstlisting}

\subsection{Algebraic Data Types}

\begin{lstlisting}
data Shape = Circle Number
           | Rectangle Number Number
           | Point

area :: Shape -> Number
area (Circle r)      = 3.14159 * r * r
area (Rectangle w h) = w * h
area Point           = 0.0
\end{lstlisting}

\subsection{Records}

Named fields with accessor functions:

\begin{lstlisting}
data Person = Person { name :: String
                     , age  :: Integer }

greet :: Person -> String
greet p = "Hello, " <> name p <> "!"

-- Record construction
alice :: Person
alice = Person { name = "Alice", age = 30 }

-- Record update (creates a copy)
olderAlice :: Person
olderAlice = alice { age = 31 }
\end{lstlisting}

\subsection{Newtypes}

Zero-cost wrappers:

\begin{lstlisting}
newtype Metres = Metres Number
newtype UserId = UserId Integer
\end{lstlisting}

\subsection{Type Aliases}

\begin{lstlisting}
type Pair a = (a, a)
type Name = String
\end{lstlisting}

\subsection{Maybe and Either}

\begin{lstlisting}
safeDivide :: Integer -> Integer -> Maybe Integer
safeDivide _ 0 = Nothing
safeDivide a b = Just (a `div` b)

parseAge :: String -> Either String Integer
parseAge s = let n = read s
             in if n >= 0 then Right n
                else Left "negative age"
\end{lstlisting}

%======================================================================
\section{Functions}

\subsection{Definition and Application}

\begin{lstlisting}
add :: Integer -> Integer -> Integer
add a b = a + b

-- Application (no parentheses needed)
result :: Integer
result = add 3 4    -- 7

-- Partial application
inc :: Integer -> Integer
inc = add 1
\end{lstlisting}

\subsection{Lambda Expressions}

\begin{lstlisting}
double :: Integer -> Integer
double = \x -> x * 2

-- Multi-parameter
add :: Integer -> Integer -> Integer
add = \a b -> a + b

-- Pattern matching in lambdas
fstOf :: (a, b) -> a
fstOf = \(x, _) -> x
\end{lstlisting}

\subsection{Operators}

\begin{lstlisting}
-- Operators are binary functions
result = 1 + 2 * 3

-- Prefix form with parentheses
result = (+) 1 2

-- Functions as infix with backticks
result = 10 `div` 3

-- Sections
double = (* 2)
inc    = (+ 1)
halve  = (`div` 2)
\end{lstlisting}

\subsubsection{Operator Precedence}

Declare fixity with \texttt{infixl}, \texttt{infixr}, or \texttt{infix}
(levels 0--9):

\begin{lstlisting}
infixl 6 +
infixl 6 -
infixl 7 *
infixr 5 :
infixr 0 $
\end{lstlisting}

Standard Haskell fixities are the defaults.

\subsection{Guards}

\begin{lstlisting}
abs :: Integer -> Integer
abs n | n < 0     = -n
      | otherwise = n

classify :: Integer -> String
classify n
    | n == 0    = "zero"
    | n > 0     = "positive"
    | otherwise = "negative"
\end{lstlisting}

\subsection{Where Clauses}

\begin{lstlisting}
circleArea :: Number -> Number
circleArea r = pi * r * r
    where pi = 3.14159265

bmi :: Number -> Number -> String
bmi weight height
    | b < 18.5  = "underweight"
    | b < 25.0  = "normal"
    | otherwise = "overweight"
    where b = weight / (height * height)
\end{lstlisting}

\subsection{Let/In Expressions}

\begin{lstlisting}
cylinderVolume :: Number -> Number -> Number
cylinderVolume r h =
    let pi = 3.14159265
        baseArea = pi * r * r
    in baseArea * h
\end{lstlisting}

%======================================================================
\section{Pattern Matching}

\subsection{Function Definitions}

\begin{lstlisting}
factorial :: Integer -> Integer
factorial 0 = 1
factorial n = n * factorial (n - 1)

head :: [a] -> a
head (x:_) = x

fst :: (a, b) -> a
fst (a, _) = a
\end{lstlisting}

\subsection{Case Expressions}

\begin{lstlisting}
describe :: [a] -> String
describe xs = case xs of
    []    -> "empty"
    [_]   -> "singleton"
    _:_:_ -> "multiple"
\end{lstlisting}

\subsection{Guards in Case}

\begin{lstlisting}
eval :: Expr -> Integer
eval e = case e of
    Lit n      -> n
    Add a b    -> eval a + eval b
    Neg n | n < 0 -> eval (Neg (Lit (-n)))
          | otherwise -> -(eval n)
\end{lstlisting}

%======================================================================
\section{Typeclasses}

\subsection{Defining a Class}

\begin{lstlisting}
class Printable a where
    display :: a -> String
\end{lstlisting}

\subsection{Instances}

\begin{lstlisting}
instance Printable Integer where
    display n = show n

instance Printable Bool where
    display True  = "yes"
    display False = "no"
\end{lstlisting}

\subsection{Superclass Constraints}

\begin{lstlisting}
instance Show a => Printable [a] where
    display xs = "[" <> intercalate ", " (map show xs) <> "]"
\end{lstlisting}

\subsection{Deriving}

Automatic instance generation for common classes:

\begin{lstlisting}
data Color = Red | Green | Blue
    deriving (Show, Eq, Ord, Enum, Bounded)

data Tree a = Leaf a | Branch (Tree a) (Tree a)
    deriving (Show, Eq, Functor)
\end{lstlisting}

Supported: \texttt{Show}, \texttt{Eq}, \texttt{Ord}, \texttt{Enum},
\texttt{Bounded}, \texttt{Functor}, and \texttt{LuaDict}.
\texttt{LuaDict} is special: it generates no methods but changes the
record's runtime representation for Lua interop --- see
Section~\ref{sec:luadict}.

\subsection{Built-in Classes}

\begin{tabular}{@{}ll@{}}
\toprule
\textbf{Class} & \textbf{Methods} \\
\midrule
\texttt{Show}        & \texttt{show :: a -> String} \\
\texttt{Eq}          & \texttt{(==), (/=) :: a -> a -> Bool} \\
\texttt{Ord}         & \texttt{compare, (<), (>), (<=), (>=)} \\
\texttt{Enum}        & \texttt{succ, pred, toEnum, fromEnum, enumFrom, \ldots} \\
\texttt{Bounded}     & \texttt{minBound, maxBound} \\
\texttt{Read}        & \texttt{read :: String -> a} \\
\texttt{Functor}     & \texttt{fmap :: (a -> b) -> f a -> f b} \\
\texttt{Applicative} & \texttt{pure, (<*>)} \\
\texttt{Monad}       & \texttt{(>>=), return} \\
\bottomrule
\end{tabular}

%======================================================================
\section{Monads and IO}

\subsection{IO Actions}

Programs with a \texttt{main :: IO ()} are standalone executables:

\begin{lstlisting}
main :: IO ()
main = do
    putStrLn "What is your name?"
    name <- getLine
    putStrLn ("Hello, " <> name <> "!")
\end{lstlisting}

\subsection{Do-Notation}

Do-notation desugars to monadic bind:

\begin{lstlisting}
-- These are equivalent:
example1 :: IO ()
example1 = do
    x <- getLine
    let y = x <> "!"
    putStrLn y

example2 :: IO ()
example2 = getLine >>= \x ->
           let y = x <> "!" in
           putStrLn y
\end{lstlisting}

\subsection{Common IO Functions}

\begin{lstlisting}
putStrLn :: String -> IO ()     -- print with newline
putStr   :: String -> IO ()     -- print without newline
getLine  :: IO String           -- read a line
getArgs  :: IO [String]         -- command-line arguments
exit     :: ExitValue -> IO ()  -- exit process
error    :: String -> a         -- abort with message
\end{lstlisting}

\subsection{Exception Handling}

\begin{lstlisting}
main :: IO ()
main = do
    result <- try (readFile "data.txt")
    case result of
        Right contents -> putStrLn contents
        Left err       -> putStrLn ("Error: " <> err)

-- Or with catch:
safeRead :: String -> IO String
safeRead path = catch (readFile path)
                      (\err -> return "")
\end{lstlisting}

%======================================================================
\section{Evaluation Strategy}

MATA-LL is non-strict: function arguments are not evaluated until needed.
Values are memoized once forced.

\begin{lstlisting}
-- Infinite list -- only computed as needed
nats :: [Integer]
nats = [1..]

fibs :: [Integer]
fibs = 1 : 1 : zipWith (+) fibs (tail fibs)

-- take 10 fibs evaluates only 10 elements
first10 :: [Integer]
first10 = take 10 fibs
\end{lstlisting}

\subsection{Explicit Forcing}

Use \texttt{seq} to force evaluation:

\begin{lstlisting}
seq :: a -> b -> b
-- Forces first argument, returns second

strictSum :: [Integer] -> Integer
strictSum = go 0
    where go acc []     = acc
          go acc (x:xs) = seq acc (go (acc + x) xs)
\end{lstlisting}

\begin{note}
The compiler performs demand analysis and cheapness analysis automatically.
Cheap expressions (literals, variables, simple arithmetic) are never
thunked.  Most programs do not need manual \texttt{seq}.
\end{note}

%======================================================================
\section{List Comprehensions}

\begin{lstlisting}
-- Generators and guards
evens :: [Integer]
evens = [x | x <- [1..20], x `mod` 2 == 0]

-- Multiple generators (nested loops)
pairs :: [(Integer, Integer)]
pairs = [(x, y) | x <- [1..3], y <- [1..3], x /= y]

-- Pattern-matching generators
successes :: [Result a] -> [a]
successes results = [v | Ok v <- results]
\end{lstlisting}

%======================================================================
\section{Modules}

\subsection{Importing}

\begin{lstlisting}
-- Import everything from a module
import Data.List

-- Import specific items
import Data.List (sortBy, groupBy)

-- Import hiding specific items
import Data.List hiding (head, tail)

-- Qualified import
import qualified Data.Map as M
\end{lstlisting}

Modules map to file paths: \texttt{Data.List} resolves to
\texttt{Data/List.mll} relative to source or library directories.

\subsection{Available Library Modules}

\begin{tabular}{@{}ll@{}}
\toprule
\textbf{Import} & \textbf{Provides} \\
\midrule
\texttt{ByteString}      & Byte sequence operations \\
\texttt{LIO}             & File I/O (\texttt{fOpen}, \texttt{fRead}, \texttt{fWrite}, \ldots) \\
\texttt{LMath}           & \texttt{math.*} bindings (sin, cos, random, \ldots) \\
\texttt{LOS}             & OS functions \\
\texttt{LString}         & String utilities \\
\texttt{LBit}            & Bitwise operations \\
\texttt{Regex}           & CPS-based regex matching \\
\texttt{JSON}            & JSON parsing and generation \\
\texttt{Data.List}       & sortBy, nubBy, groupBy, intercalate, \ldots \\
\texttt{Data.Maybe}      & fromMaybe, catMaybes, mapMaybe, \ldots \\
\texttt{Data.Map}        & HashMap operations \\
\texttt{Control.Monad}   & mapM, forM, sequence, \ldots \\
\bottomrule
\end{tabular}

The Prelude is imported automatically into every module.

%======================================================================
\section{Foreign Function Interface}

MATA-LL interoperates with Lua through type families.

\subsection{Pure FFI}

For Lua functions with no side effects:

\begin{lstlisting}
sin :: Number -> LuaPure "math.sin" Number
cos :: Number -> LuaPure "math.cos" Number
strlen :: String -> LuaPure "string.len" Integer
\end{lstlisting}

\texttt{LuaPure "name" T} reduces to \texttt{T} --- the function is
callable as a pure expression.

\subsection{Effectful FFI}

For Lua functions that perform I/O or have side effects:

\begin{lstlisting}
rnd :: Integer -> Integer -> LuaIO "math.random" Integer
clock :: LuaIO "os.clock" Number
\end{lstlisting}

\texttt{LuaIO "name" T} reduces to \texttt{IO T}.

\subsection{Method-Call Syntax}

Colon prefix generates Lua method calls (\texttt{obj:method()}):

\begin{lstlisting}
fileWrite :: Handle -> String -> LuaIO ":write" ()
fileClose :: Handle -> LuaIO ":close" ()
\end{lstlisting}

\subsection{Optional Parameters}

\texttt{Maybe} arguments become optional Lua parameters:

\begin{lstlisting}
rnd :: Number -> Maybe Number -> LuaIO "math.random" Number
-- rnd 1.0 Nothing     => math.random(1.0)
-- rnd 1.0 (Just 6.0)  => math.random(1.0, 6.0)
\end{lstlisting}

\subsection{Iterators}

Wrap Lua iterator protocols:

\begin{lstlisting}
gmatch :: String -> String -> LuaIterator "string.gmatch" String
\end{lstlisting}

\subsection{Exporting to Lua}

Use \texttt{export} to make functions callable from Lua:

\begin{lstlisting}
fib :: [Integer]
fib = 1 : 1 : zipWith (+) fib (tail fib)

export fibonacci :: Integer -> [Integer]
fibonacci = flip take fib
\end{lstlisting}

From Lua:
\begin{lstlisting}[language={}]
local mll = require("fib")
local fibs = mll.fibonacci(10)
-- fibs is a plain Lua array: {1, 1, 2, 3, 5, 8, 13, 21, 34, 55}
\end{lstlisting}

Exported functions automatically convert MLL values to Lua values
(forces thunks, converts cons-lists to arrays).

\subsection{LuaDict Records}
\label{sec:luadict}

Many Lua APIs take or return \emph{dictionaries} --- tables keyed by
name, such as \texttt{\{width = 80, height = 25\}}.  Deriving
\texttt{LuaDict} on a record makes its runtime representation exactly
such a table.  This is not a marshalling layer that converts at the FFI
boundary: the record value \emph{is} a name-keyed Lua table from the
moment it is constructed, so it can be handed to (or received from) Lua
functions with no conversion cost.

\begin{lstlisting}
data Config = Config { width  :: Integer
                     , height :: Integer
                     , title  :: String
                     } deriving (Show, Eq, LuaDict)

-- a Lua function expecting a dictionary
setup :: Config -> LuaIO "display.setup" ()

main :: IO ()
main = do
    let c = Config { width = 80, height = 25, title = "hi" }
    -- at runtime, c is {width = 80, height = 25, title = "hi"}
    setup c
\end{lstlisting}

Without \texttt{LuaDict}, records are stored as positional Lua arrays;
with it, every field becomes a named key.  On the MATA-LL side nothing
changes: accessor functions, record construction (named or positional),
record update, pattern matching, and derived \texttt{Show}/\texttt{Eq}
all work as usual.  The direction is symmetric --- an FFI function
declared to \emph{return} a \texttt{LuaDict} type accepts a Lua table
with the corresponding keys:

\begin{lstlisting}
currentConfig :: LuaIO "display.get_config" Config
\end{lstlisting}

\texttt{LuaDict} places requirements on the data type, because a
name-keyed table has no constructor tag and needs field names for its
keys.  The type must have:

\begin{itemize}[nosep]
  \item exactly one constructor,
  \item record syntax (named fields, not positional),
  \item at least one field,
  \item no GADT or existential constructor.
\end{itemize}

Violations are compile errors with an explanation of why.

\subsubsection{Renaming Keys with \texttt{as}}

By default the Lua key is the field name.  Since record accessors share
a single global function namespace (as in Haskell), fields are
conventionally prefixed (\texttt{cfgWidth}, \texttt{cfgHeight}) --- but
a Lua API will expect \texttt{width}, not \texttt{cfgWidth}.  The
\texttt{as} annotation renames the Lua table key:

\begin{lstlisting}
data Config = Config { cfgWidth  as "width"  :: Integer
                     , cfgHeight as "height" :: Integer
                     , cfgTitle  as "title"  :: String
                     } deriving (Show, Eq, LuaDict)

area :: Config -> Integer
area c = cfgWidth c * cfgHeight c   -- Haskell name, unchanged

resize :: Config -> Integer -> Config
resize c w = c { cfgWidth = w }     -- updates the "width" key
\end{lstlisting}

\texttt{as "key"} changes \emph{only} the key in the runtime Lua table.
In MATA-LL source the field keeps its declared name everywhere:
accessors, record construction, record update, and pattern matching all
use \texttt{cfgWidth}.  The rename applies in both directions across
the FFI boundary --- tables you pass to Lua carry the renamed keys, and
tables returned from Lua are read by the renamed keys.

Besides dropping prefixes, \texttt{as} is useful for
\texttt{snake\_case} keys expected by Lua APIs, and for keys that are
Lua reserved words:

\begin{lstlisting}
data Creds = Creds { credsUser as "user_name" :: String
                   , credsNote as "function"  :: String
                   } deriving (LuaDict)
\end{lstlisting}

Reserved-word keys are emitted with bracket syntax
(\texttt{\{["function"] = ...\}}), so any string is a valid key.

The compiler enforces two rules on renames:

\begin{itemize}[nosep]
  \item Effective keys must be unique and non-empty.  Each field
        becomes one key of a single table, so two fields mapping to the
        same key (or a key of \texttt{""}) are rejected.
  \item \texttt{as} requires \texttt{deriving (LuaDict)}.  Without it
        the record is a positional array, there is no key to rename,
        and the annotation is rejected rather than silently ignored.
\end{itemize}

\begin{note}
\texttt{as} field renaming is a MATA-LL extension with no GHC
equivalent.  In GHC, \texttt{as} is only a keyword in import lists;
here it is additionally recognized inside record declarations.
\end{note}

%======================================================================
\section{Advanced Types}

\subsection{GADTs}

Generalized Algebraic Data Types allow constructors to refine type
parameters:

\begin{lstlisting}
data Expr a where
    LitI :: Integer -> Expr Integer
    LitB :: Bool -> Expr Bool
    Add  :: Expr Integer -> Expr Integer -> Expr Integer
    If   :: Expr Bool -> Expr a -> Expr a -> Expr a

eval :: Expr a -> a
eval (LitI n)     = n
eval (LitB b)     = b
eval (Add x y)    = eval x + eval y
eval (If c t e)   = if eval c then eval t else eval e
\end{lstlisting}

\subsection{Existential Types}

Hide a type behind an interface:

\begin{lstlisting}
data ShowBox = forall a. MkShowBox a (a -> String)

showIt :: ShowBox -> String
showIt (MkShowBox x f) = f x

boxes :: [ShowBox]
boxes = [MkShowBox 42 show, MkShowBox "hi" id]
\end{lstlisting}

\subsection{DataKinds}

Promote data constructors to the type level for compile-time safety:

\begin{lstlisting}
data AgendaState = Empty | NonEmpty

data Agenda a (s :: AgendaState) where
    Nil   :: Agenda a 'Empty
    One   :: a -> Agenda a 'NonEmpty
    Cons  :: a -> Agenda a 'NonEmpty -> Agenda a 'NonEmpty

-- Only non-empty agendas can be submitted
submit :: Agenda String 'NonEmpty -> IO ()
submit agenda = putStrLn "Submitted!"
\end{lstlisting}

\subsection{Type Families}

User-defined type-level computation:

\begin{lstlisting}
type family Element container where
    Element [a]           = a
    Element (HashMap k v) = v
\end{lstlisting}

\subsection{Rank-2 Types and Scoped Mutation}

The \texttt{ST} monad provides mutable arrays inside pure computations:

\begin{lstlisting}
import Data.List (replicate)

sumList :: [Integer] -> Integer
sumList xs = runST (do
    acc <- newSTArrayFromList [0]
    go acc xs
    readSTArray acc 0)

go :: STArray s -> [Integer] -> ST s ()
go acc []     = return ()
go acc (x:rest) = do
    cur <- readSTArray acc 0
    writeSTArray acc 0 (cur + x)
    go acc rest
\end{lstlisting}

The \texttt{forall s.} in \texttt{runST}'s type prevents mutable state
from escaping.

\subsection{Scoped Lua Callbacks}

When exporting functions that receive Lua callbacks:

\begin{lstlisting}
export processEvent :: forall s. (Integer -> Integer -> LuaIO s Integer)
                    -> Integer -> LuaIO s Integer
processEvent f n = do
    (liftIO . putStrLn) $ "Called with " <> show n
    f n (n + 1)
\end{lstlisting}

The \texttt{forall s.} scope parameter prevents the callback from being
stored or used outside its invocation.

%======================================================================
\section{HashMap}

A built-in dictionary type backed by Lua tables:

\begin{lstlisting}
import Data.Map

main :: IO ()
main = do
    let m = hashmap_fromList [("a", 1), ("b", 2), ("c", 3)]
    putStrLn (show (hashmap_lookup "b" m))  -- Just 2
    let m2 = hashmap_insert "d" 4 m
    putStrLn (show (hashmap_keys m2))
    putStrLn (show (hashmap_member "x" m))  -- False
\end{lstlisting}

Operations: \texttt{hashmap\_lookup}, \texttt{hashmap\_insert},
\texttt{hashmap\_delete}, \texttt{hashmap\_keys}, \texttt{hashmap\_values},
\texttt{hashmap\_member}, \texttt{hashmap\_fromList}.

%======================================================================
\section{Ranges and Enumerations}

Range syntax desugars to \texttt{Enum} methods:

\begin{lstlisting}
[1..10]       -- enumFromTo 1 10       => [1,2,...,10]
[1,3..10]     -- enumFromThenTo 1 3 10 => [1,3,5,7,9]
[1..]         -- enumFrom 1            => [1,2,3,...] (infinite)
[1,3..]       -- enumFromThen 1 3      => [1,3,5,...] (infinite)
\end{lstlisting}

Works with any type that has an \texttt{Enum} instance, including
user-defined types with \texttt{deriving Enum}.

%======================================================================
\section{Differences from Haskell}

If you are coming from Haskell, keep these differences in mind:

\begin{enumerate}[nosep]
  \item \textbf{String is not \texttt{[Char]}}.  It is a native Lua string.
        Use \texttt{<>} for concatenation; \texttt{++} is list-append only.
  \item \textbf{No \texttt{Num} class}.  Arithmetic operators work directly
        on \texttt{Integer} and \texttt{Number}.  Numeric literals are not
        polymorphic.
  \item \textbf{No \texttt{Char} type}.  Character operations use string
        functions from \texttt{LString}.
  \item \textbf{No multi-parameter typeclasses}.
  \item \textbf{No lazy I/O}.  File operations read eagerly.
  \item \textbf{No \texttt{IORef}/\texttt{MVar}/\texttt{STRef}}.
        Use \texttt{STArray} for scoped mutation.
  \item \textbf{All top-level bindings need type signatures}.
  \item \textbf{Single-file modules}.  Each \texttt{.mll} file is one module.
\end{enumerate}

%======================================================================
\section{Complete Example: AVL Dictionary}

\begin{lstlisting}
data Dict k v = Empty
              | Node Integer k v (Dict k v) (Dict k v)

empty :: Dict k v
empty = Empty

height :: Dict k v -> Integer
height Empty = 0
height (Node h _ _ _ _) = h

node :: k -> v -> Dict k v -> Dict k v -> Dict k v
node k v l r = Node (1 + max (height l) (height r)) k v l r

insert :: k -> v -> Dict k v -> Dict k v
insert k v Empty = node k v Empty Empty
insert k v (Node h nk nv left right)
    | k < nk    = rebalance (node nk nv (insert k v left) right)
    | k > nk    = rebalance (node nk nv left (insert k v right))
    | otherwise = Node h k v left right

lookup' :: k -> Dict k v -> Maybe v
lookup' _ Empty = Nothing
lookup' k (Node _ nk nv left right)
    | k < nk    = lookup' k left
    | k > nk    = lookup' k right
    | otherwise = Just nv

main :: IO ()
main = do
    let d = insert "c" 3 $ insert "b" 2 $ insert "a" 1 $ empty
    putStrLn (show (lookup' "b" d))   -- Just 2
    putStrLn (show (lookup' "z" d))   -- Nothing
\end{lstlisting}

(Full rebalancing omitted for brevity --- see \texttt{examples/dict.mll}.)

%======================================================================
\section{Complete Example: Lazy Fibonacci Export}

A library module that exports an efficient Fibonacci function to Lua:

\begin{lstlisting}
fib :: [Integer]
fib = 1 : 1 : zipWith (+) fib (tail fib)

export fibonacci :: Integer -> [Integer]
fibonacci = flip take fib
\end{lstlisting}

Usage from Lua:

\begin{lstlisting}[language={}]
local m = require("fib")
for i, v in ipairs(m.fibonacci(20)) do
    print(i, v)
end
\end{lstlisting}

The infinite list is computed lazily; only the requested elements are
evaluated.

%======================================================================
\section{Prelude Reference}

\subsection{Functions}

\begin{lstlisting}
-- Output
putStrLn :: String -> IO ()
putStr   :: String -> IO ()

-- Conversion
show :: Show a => a -> String
read :: Read a => String -> a

-- Combinators
id    :: a -> a
const :: a -> b -> a
flip  :: (a -> b -> c) -> b -> a -> c
(.)   :: (b -> c) -> (a -> b) -> a -> c
($)   :: (a -> b) -> a -> b

-- Arithmetic (work on Integer and Number)
(+), (-), (*) :: ...
div, mod      :: Integer -> Integer -> Integer
(/)           :: Number -> Number -> Number

-- Comparison (require Eq/Ord instances)
(==), (/=) :: Eq a  => a -> a -> Bool
(<), (>), (<=), (>=) :: Ord a => a -> a -> Bool
max, min   :: Ord a => a -> a -> a
compare    :: Ord a => a -> a -> Ordering

-- Boolean
(&&), (||) :: Bool -> Bool -> Bool
not        :: Bool -> Bool
otherwise  :: Bool    -- = True

-- List operations
map       :: (a -> b) -> [a] -> [b]
filter    :: (a -> Bool) -> [a] -> [a]
foldl     :: (b -> a -> b) -> b -> [a] -> b
foldr     :: (a -> b -> b) -> b -> [a] -> b
head      :: [a] -> a
tail      :: [a] -> [a]
take      :: Integer -> [a] -> [a]
drop      :: Integer -> [a] -> [a]
length    :: [a] -> Integer
reverse   :: [a] -> [a]
elem      :: Eq a => a -> [a] -> Bool
zipWith   :: (a -> b -> c) -> [a] -> [b] -> [c]
concatMap :: (a -> [b]) -> [a] -> [b]
(++)      :: [a] -> [a] -> [a]

-- String
(<>)      :: String -> String -> String

-- Tuple
fst :: (a, b) -> a
snd :: (a, b) -> b

-- Control
seq       :: a -> b -> b
error     :: String -> a
undefined :: a

-- Monadic
when   :: Bool -> IO () -> IO ()
mapM_  :: (a -> IO ()) -> [a] -> IO ()
\end{lstlisting}

\subsection{Types}

\begin{lstlisting}
data Maybe a = Nothing | Just a
data Either a b = Left a | Right b
data Ordering = LT | EQ | GT
data ExitValue = Normal | Err Integer
\end{lstlisting}

\end{document}
